\documentclass[twoside,11pt]{article}

% JMLR/MLOSS submissions normally use jmlr2e.sty from the JMLR author kit.
% Compile with the JMLR style file available on the TeX path.
\usepackage{jmlr2e}
\usepackage{booktabs}
\usepackage{url}
\usepackage{amsmath}

\jmlrheading{27}{2026}{1-}{TODO}{TODO}{TODO}{TODO}
\ShortHeadings{GraphSlim: A PyTorch Library for Graph Reduction}{Gong, Ni, Sachdeva, Yang, and Jin}
\firstpageno{1}

\begin{document}

\title{GraphSlim: A PyTorch Library for Graph Reduction}

\author{
\name Shengbo Gong \email TODO \\
\addr TODO
\AND
\name Juntong Ni \email TODO \\
\addr TODO
\AND
\name Noveen Sachdeva \email TODO \\
\addr TODO
\AND
\name Carl Yang \email TODO \\
\addr TODO
\AND
\name Wei Jin \email TODO \\
\addr TODO
}

\editor{TODO}

\maketitle

\begin{abstract}
Graph reduction methods aim to replace a large graph with a smaller graph that preserves
downstream predictive performance or structural properties. We introduce GraphSlim, an
open-source PyTorch and PyTorch Geometric library for graph reduction. GraphSlim provides
a unified implementation of representative methods from the three main families of graph
reduction: sparsification, coarsening, and condensation. The package includes dataset
loading and conversion utilities, reduction method construction, downstream evaluation,
benchmark scripts, optional experiment tracking, and local visualization helpers. By
standardizing the interfaces between datasets, reduction methods, and evaluators, GraphSlim
lowers the engineering barrier for comparing graph reduction algorithms and for applying
them to graph learning workflows.
\end{abstract}

\begin{keywords}
graph reduction, graph condensation, graph sparsification, graph coarsening, graph neural
networks, PyTorch, PyTorch Geometric, open-source software
\end{keywords}

\section{Introduction}

Graph neural networks have become a standard tool for learning from relational data, but
training and evaluating models on large graphs can be expensive in memory, time, and
storage. Graph reduction addresses this problem by constructing a smaller graph or
training set that retains useful information from the original graph. The resulting reduced
object can accelerate downstream model selection, repeated training, privacy-sensitive
sharing, visualization, and controlled benchmarking.

GraphSlim is a Python package for graph reduction research and practice. It is designed
around a PyTorch-first workflow and accepts graph data in PyTorch Geometric-style formats.
The package exposes implementations of representative algorithms across sparsification,
coarsening, and condensation, together with common evaluation and benchmarking utilities.
GraphSlim is distributed as open-source software and is available from PyPI and GitHub.

The main design goal of GraphSlim is to make graph reduction methods comparable under a
shared data and evaluation contract. Instead of each method carrying its own ad hoc
orchestration logic, recent versions organize method construction through a shared registry
and factory. This allows entry points such as command-line training, initialization, and
package-level usage to construct reduction methods consistently.

\section{Related Software and Motivation}

The Python graph learning ecosystem includes mature libraries such as PyTorch
Geometric~\citep{fey2019fast}, DGL~\citep{wang2019dgl}, NetworkX~\citep{hagberg2008networkx},
and OGB~\citep{hu2020ogb}. These packages provide graph data structures, message passing
operators, datasets, and evaluation protocols, but they are not primarily graph reduction
toolkits. GraphSlim complements these libraries by focusing on the algorithmic layer that
constructs reduced graphs and reduced training sets.

Graph reduction has several distinct methodological traditions. Sparsification methods
select nodes or edges, coarsening methods contract or aggregate graph structure, and
condensation methods synthesize features, labels, and sometimes graph structure. A recent
survey by \citet{hashemi2024comprehensive} and the GC4NC benchmark of
\citet{gong2025gc4nc} motivate a unified implementation layer for these families.
GraphSlim follows this motivation by providing a package where multiple reduction families
share dataset loaders, evaluation routines, and experiment configuration.

\section{Software Overview}

GraphSlim is organized into modules that separate datasets, methods, models, evaluation,
tracking, and visualization. Table~\ref{tab:modules} summarizes the main package
components.

\begin{table}[h]
\centering
\begin{tabular}{ll}
\toprule
Module & Role \\
\midrule
\texttt{graphslim.dataset} & Dataset loading, splitting, conversion, attacks, and persistence \\
\texttt{graphslim.sparsification} & Node, core-set, and edge sparsification methods \\
\texttt{graphslim.coarsening} & Graph coarsening and contraction methods \\
\texttt{graphslim.condensation} & Synthetic graph and feature condensation methods \\
\texttt{graphslim.models} & GNN backbones, graph layers, kernels, and parametrized adjacency modules \\
\texttt{graphslim.evaluation} & Downstream evaluation, timing, memory, NAS, and graph property evaluation \\
\texttt{graphslim.reduction} & Method registry and factory for constructing reducers \\
\texttt{graphslim.tracking} & Optional experiment tracking and graph summary logging \\
\texttt{graphslim.visualization} & Local visualization utilities for small graph comparisons \\
\bottomrule
\end{tabular}
\caption{Main GraphSlim modules and responsibilities.}
\label{tab:modules}
\end{table}

The package currently supports representative methods from all three reduction families.
The method registry includes node and edge sparsification methods such as \texttt{kcenter},
\texttt{herding}, \texttt{random}, \texttt{random\_edge}, \texttt{g\_spar}, and
\texttt{t\_spanner}; coarsening methods such as \texttt{variation\_neighborhoods},
\texttt{heavy\_edge}, \texttt{kron}, and \texttt{clustering}; and condensation methods such
as \texttt{gcond}, \texttt{doscond}, \texttt{sfgc}, \texttt{msgc}, \texttt{sgdd},
\texttt{gcsntk}, \texttt{geom}, \texttt{simgc}, \texttt{gcdm}, and \texttt{gecc}.

\section{Usage}

GraphSlim can be installed from PyPI or from source. The command-line interface exposes a
single training entry point for loading a dataset, running a reduction method, and evaluating
the reduced graph. A typical command is:

\begin{verbatim}
graphslim-train --dataset cora --method kcenter --gpu_id -1 \
  --epochs 100 --eval_epochs 300 --reduction_rate 0.5
\end{verbatim}

The package can also be used directly from Python. The following sketch shows the intended
programmatic workflow:

\begin{verbatim}
from graphslim.config import cli
from graphslim.dataset import get_dataset
from graphslim.evaluation import Evaluator
from graphslim.reduction import create_reducer

args = cli(standalone_mode=False)
args.dataset = "cora"
args.method = "gcond"
args.reduction_rate = 0.5

graph = get_dataset(args.dataset, args=args, load_path=args.load_path)
reducer = create_reducer(args.method, setting=args.setting, data=graph, args=args)
reduced_graph = reducer.reduce(graph, verbose=True)

evaluator = Evaluator(args)
mean, std = evaluator.evaluate(reduced_graph, model_type=args.eval_model)
\end{verbatim}

This interface is intentionally method-agnostic: changing the reduction method should not
require rewriting the dataset loading or evaluation code. Method-specific hyperparameters
are handled through JSON configuration files and command-line options.

\section{Datasets and Evaluation}

GraphSlim wraps several dataset sources behind a common \texttt{TransAndInd} data object.
The current loaders cover Planetoid datasets, PyTorch Geometric datasets such as Flickr,
Reddit2, Coauthor, CitationFull, and Amazon, OGB node property prediction datasets, a
GraphSAINT-style loader for \texttt{ogbn-arxiv}, and DGL fraud datasets converted to a PyG
style interface.

The evaluator trains downstream GNN models on the reduced graph and reports validation and
test metrics. Supported evaluation backbones include GCN, SGC, APPNP, Cheby, GraphSage,
GAT, and SGFormer. The evaluation layer is shared across reduction families, which makes it
possible to compare sparsification, coarsening, and condensation methods under the same
metric and split definitions.

\section{Implementation Notes}

The shared registry in \texttt{graphslim.reduction} maps public method names to concrete
implementation classes and imports them lazily. Lazy importing keeps package import time
low for users who only need method discovery or metadata. The same registry is used by the
command-line training entry point and by condensation initialization routines, reducing
duplicated method-selection logic.

GraphSlim keeps dataset conversion at package boundaries. PyG-style tensors and sparse
matrices form the internal data contract, while compatibility utilities support conversion
from DGL and related graph objects. Optional tracking and visualization utilities are kept
separate from the core training path so that users can install only the dependencies needed
for their workflow.

\section{Quality Control and Reproducibility}

GraphSlim includes unit tests for compatibility helpers, the reduction registry, tracking,
and visualization. Recent smoke tests exercise the public training path across all registered
methods on Citeseer with minimum CPU settings, and a representative module-by-loader matrix
covers Planetoid, PyG Amazon, PyG Coauthor, PyG Flickr, GraphSAINT/OGB, and DGL
FraudDataset loading paths.

These smoke tests are intended to check pipeline health rather than benchmark model
quality. In particular, very short evaluation runs can produce repeated deterministic
accuracy values and should not be interpreted as scientific performance comparisons.
Quality benchmarking should use method-appropriate reduction epochs, enough downstream
evaluation epochs, multiple random seeds, and the benchmark scripts distributed with the
project.

\section{Availability}

GraphSlim is open-source software released under the MIT license. The package is available
on PyPI as \texttt{graphslim}, with source code, documentation, examples, benchmark scripts,
and a web visualization interface linked from the project repository:

\begin{center}
\url{https://github.com/Emory-Melody/GraphSlim}
\end{center}

The package supports Python 3.9--3.11 and depends on PyTorch, PyTorch Geometric, SciPy,
NumPy, scikit-learn, Networkit, PyGSP, OGB, and related graph learning utilities.

\section{Limitations and Future Work}

GraphSlim is a research-oriented package and some methods inherit the computational cost
and dependency complexity of their original algorithms. Condensation methods that rely on
teacher trajectory buffers, such as \texttt{sfgc} and \texttt{geom}, can be expensive in
small smoke-test settings unless replay buffers are precomputed or lightweight test buffers
are used. Some coarsening implementations depend on external graph libraries whose APIs
may change across versions, and some large-dataset experiments require substantial memory
or preprocessing time.

Future work includes expanding fast CPU tests for method construction and reduction,
strengthening cross-dataset CI coverage, adding optional compiled acceleration for repeated
sparse conversion and contraction routines, and continuing to align the public API with
the unified registry and evaluation pipeline.

\section{Conclusion}

GraphSlim provides a unified PyTorch-based library for graph reduction. By collecting
sparsification, coarsening, and condensation methods behind shared dataset and evaluation
interfaces, it supports reproducible comparison and lowers the cost of using graph reduction
in graph learning workflows.

\acks{TODO: Add funding, contributor, and infrastructure acknowledgments.}

\begin{thebibliography}{}

\bibitem[Fey and Lenssen(2019)]{fey2019fast}
Fey, M. and Lenssen, J. E. (2019).
\newblock Fast graph representation learning with PyTorch Geometric.
\newblock In \emph{ICLR Workshop on Representation Learning on Graphs and Manifolds}.

\bibitem[Gong et~al.(2025)]{gong2025gc4nc}
Gong, S., Ni, J., Sachdeva, N., Yang, C., and Jin, W. (2025).
\newblock {GC}4{NC}: A benchmark framework for graph condensation on node
classification with new insights.
\newblock In \emph{Advances in Neural Information Processing Systems, Datasets and
Benchmarks Track}.

\bibitem[Hagberg et~al.(2008)]{hagberg2008networkx}
Hagberg, A. A., Schult, D. A., and Swart, P. J. (2008).
\newblock Exploring network structure, dynamics, and function using NetworkX.
\newblock In \emph{Proceedings of the 7th Python in Science Conference}.

\bibitem[Hashemi et~al.(2024)]{hashemi2024comprehensive}
Hashemi, M., Gong, S., Ni, J., Fan, W., Prakash, B. A., and Jin, W. (2024).
\newblock A comprehensive survey on graph reduction: Sparsification, coarsening,
and condensation.
\newblock In \emph{Proceedings of the Thirty-Third International Joint Conference on
Artificial Intelligence}, pages 8058--8066.

\bibitem[Hu et~al.(2020)]{hu2020ogb}
Hu, W., Fey, M., Zitnik, M., Dong, Y., Ren, H., Liu, B., Catasta, M., and Leskovec, J.
(2020).
\newblock Open Graph Benchmark: Datasets for machine learning on graphs.
\newblock In \emph{Advances in Neural Information Processing Systems}.

\bibitem[Paszke et~al.(2019)]{paszke2019pytorch}
Paszke, A., Gross, S., Massa, F., Lerer, A., Bradbury, J., Chanan, G., Killeen, T.,
Lin, Z., Gimelshein, N., Antiga, L., Desmaison, A., Kopf, A., Yang, E., DeVito, Z.,
Raison, M., Tejani, A., Chilamkurthy, S., Steiner, B., Fang, L., Bai, J., and Chintala,
S. (2019).
\newblock PyTorch: An imperative style, high-performance deep learning library.
\newblock In \emph{Advances in Neural Information Processing Systems}.

\bibitem[Wang et~al.(2019)]{wang2019dgl}
Wang, M., Yu, L., Zheng, D., Gan, Q., Gai, Y., Ye, Z., Li, M., Zhou, J., Huang, Q.,
Ma, C., Huang, Z., Guo, Q., Zhang, H., Lin, H., Zhao, J., Li, J., Smola, A. J., and
Zhang, Z. (2019).
\newblock Deep Graph Library: Towards efficient and scalable deep learning on graphs.
\newblock In \emph{ICLR Workshop on Representation Learning on Graphs and Manifolds}.

\end{thebibliography}

\end{document}
